\documentclass[10pt,a4paper]{amsart}
\usepackage[margin=19mm]{geometry}
\usepackage{amsmath,amssymb,mathtools}
\usepackage[scr=rsfs,cal=euler]{mathalfa}
\usepackage{xfrac}
\usepackage[unicode,hidelinks,bookmarks=false]{hyperref}
\newcommand{\ie}{i.e.\ }
\newcommand{\cf}{cf.\ }
\newcommand{\resp}{respectively\ }
\newcommand{\Subsection}{\subsection}
\providecommand{\nobreakdash}{\nobreak\hbox{-}}
\setlength{\emergencystretch}{2em}
\multlinegap0pt
\usepackage{amssymb}
\usepackage{enumerate}
\newtheorem{corollary}{Corollary}
\DeclareMathOperator{\rot}{rot}
\DeclareMathOperator{\tb}{tb}
\title{Legendrian simple knots not detected by Khovanov homology}
\author{Chun-Sheng Hsueh, Marc Kegel, David Suchodoll,, Annika Thiele}
\date{Source: arXiv:2608.24353v1 (CC BY 4.0)}
\begin{document}
\maketitle

\noindent\emph{Source and licence.} Legendrian simple knots not detected by Khovanov homology. Version arXiv:2608.24353v1 (CC BY 4.0), distributed by arXiv under the Creative Commons Attribution 4.0 International licence, http://creativecommons.org/licenses/by/4.0/, as recorded in the arXiv licence metadata for this exact version and shown on the abstract page https://arxiv.org/abs/2608.24353v1. Full licence notice: \texttt{licenses/arxiv-2608.24353-CC-BY-4.0.txt}.\par\smallskip

\noindent\textit{Editorial scope.} Counted unit: the article's corollary corollary (source0.tex lines 162--170). The article's complete original proof of it is delivered with it (source0.tex lines 172--182, one \texttt{proof} environment of 11 lines). No mathematics is written, repaired or re-derived: the statement and the proof are the article's own lines, in the article's own notation, and no retained line is substituted or re-flowed.  No further source-local context is needed: every cross-reference of every retained line resolves inside the two delivered spans.  Nothing was omitted from any retained span, so the package carries no omission record.  No retained line contains a citation, so the wrapper carries no bibliography and no bibliography entry.  No retained line contains a figure environment, an image-inclusion command or drawing code, so no image is delivered and the package declares no asset dependency.  Editorial formatting only, in addition: an AMS \texttt{amsart} preamble that restates the article's own declarations and macros used by the retained lines, namely the environments \texttt{corollary} on separate counters, together with the standard packages those lines need.  The retained environments print in this document as Corollary 1 in order of appearance, whereas the article prints the same items with the numbering of its own full document.  The complete native source of the article as distributed in the arXiv e-print for this exact version is bundled unchanged as \texttt{sources/208466/source0.tex}.
\begin{corollary}\hfill
\begin{enumerate}
    \item Any two Legendrian representatives of \(K13n1836\) obtained by stabilizing
\(L_-\) or \(L_+\) are Legendrian isotopic whenever they have the same
Thurston--Bennequin and rotation numbers.
    \item If every Legendrian representative of \(K13n1836\) with
\(\tb<-4\) destabilizes, then \(K13n1836\) is Legendrian simple.
\end{enumerate}
\end{corollary}

\begin{proof}
The representatives \(L_-\) and \(L_+\) have a common stabilization. Hence,
using that positive and negative stabilizations commute, any stabilizations of
\(L_-\) and \(L_+\) with the same classical invariants are Legendrian isotopic.
This proves~(1).

If every representative with \(\tb<-4\) destabilizes, then every Legendrian
representative destabilizes to one of the two maximal-\(tb\) representatives
\(L_-\) or \(L_+\). Thus~(1) implies that \(\tb\) and \(\rot\) determine the
Legendrian isotopy class, proving~(2).
\end{proof}

\end{document}
